\section*{Appendix C. Julia Boundary, Master Residual $\eta$, and Reversal-Probe Consolidation}
\addcontentsline{toc}{section}{Appendix C. Julia Boundary, Master Residual $\eta$, and Reversal-Probe Consolidation}

This appendix records the condensed 3.3 main line. The local residual sector is organized around the single scalar
\[
  \eta_f:=\frac13-\beta_f^2.
\]
The point of the compression is that the previously separated Hopf-gap, Julia-boundary, isotropic defect, and reversal language are different visible faces of the same transported branch quantity.

\begin{proposition}[Master residual transport and isotropic readout]
\label{lem:defect-transport}
\label{prop:master-defect-scalar}
Let $f$ range over admissible readouts of the same closure-stable kernel, and assume that on the admissible branch the visible compensated defect is isotropic,
\[
  D^{(f)}_{\mu\nu}=\eta_f\,g^{(f)}_{\mu\nu}.
\]
Assume moreover that admissible transport preserves the underlying defect class. Define
\[
  \eta_f:=\frac13-\beta_f^2.
\]
Then the following are equivalent encodings of the same local residual sector:
\begin{enumerate}[label=(\roman*)]
  \item the additive $\beta^2$-gap $\eta_f$;
  \item the continuum-limit deformation
  \[
    \delta_{\mathrm{CL},f}:=\frac{\beta_f}{1/\sqrt3}=\sqrt3\,\beta_f,
    \qquad
    \eta_f=\frac{1-\delta_{\mathrm{CL},f}^2}{3};
  \]
  \item the cosmological residuals
  \[
    \eta_f=\frac13-\Omega_M^{(f)}=\Omega_\Lambda^{(f)}-\frac23
  \]
  whenever $\Omega_M^{(f)}=\beta_f^2$ and $\Omega_\Lambda^{(f)}=1-\beta_f^2$;
  \item the isotropic defect tensor $D^{(f)}_{\mu\nu}=\eta_f g^{(f)}_{\mu\nu}$;
  \item the transported quadratic invariant
  \[
    \mathfrak I_f=\eta_f^2.
  \]
\end{enumerate}
In particular, every admissible readout carries the same local defect class through the single scalar $\eta_f$.
\end{proposition}

\begin{proof}
The displayed identities are algebraic. Once the visible defect is isotropic, every normalized quadratic contraction is determined by its scalar coefficient, hence by $\eta_f$. Transport preservation identifies the same defect class in different readouts, so the admissible residual sector is represented everywhere by the same one-scalar quantity and its powers.
\end{proof}

\begin{remark}[Symmetric value and the $10^{-3}$ scale]
For the symmetric matter/tachyon readout with $\beta_f=\beta_{\mathrm{sym}}\approx0.5493$ one gets
\[
  \eta_{\mathrm{sym}}=\frac13-\beta_{\mathrm{sym}}^2\approx 3.1603\times10^{-2},
  \qquad
  \eta_{\mathrm{sym}}^2\approx 9.987\times10^{-4}.
\]
So a number near $10^{-3}$ is most naturally read not as a primitive constant, but as the quadratic visible residual of the transported master defect scalar.
\end{remark}

\begin{proposition}[Julia boundary and exact-core principle]
\label{prop:julia-master-residual}
Assume that admissible branching does not introduce a new primitive local scalar beyond the continuum--discrete deformation of the same transport sector. Then the exact compensated core has two equivalent local descriptions:
\[
  \text{intrinsic core: }\zminus=0,
  \qquad
  \text{normalized visible core: }\eta_f=0.
\]
Moreover every admissible leading-order branch/Julia defect functional satisfies
\[
  \mathfrak J_f=c_f\,\eta_f+O(\eta_f^2)
\]
with a nonzero normalization constant $c_f$. Thus the observed Julia-boundary sensitivity is controlled, to leading order, by the same scalar $\eta_f$ that governs the isotropic defect readout.
\end{proposition}

\begin{proof}
The exact compensated core is representation-stable, but it is seen in a phase-shifted readout frame. Intrinsically the core is the null-kernel relation $\zminus=0$; after Hopf normalization of the visible branch it is the vanishing of the transported residual coefficient $\eta_f$. If admissible branching opens no new primitive scalar channel, then the first visible boundary defect must be a function of the already available scalar $\eta_f$, and its first nontrivial term is therefore linear in $\eta_f$.
\end{proof}

\begin{theorem}[Local $\eta$-completeness and restoring normal form]
\label{thm:eta-completeness-restoring}
Assume a branch-local admissible transport chart in which visible observables depend only on the local admissible branch parameter and no admissible transverse scalar channel reopens. Then:
\begin{enumerate}[label=(\roman*)]
  \item every branch-local scalar observable without derivatives is a function of $\eta$;
  \item every first-jet scalar observable is a function of $\eta$ as well, because
  \[
    \nabla_\mu \eta = A(\eta)\,U_\mu
  \]
  for some branch tangent $U$ and scalar amplitude $A$;
  \item every second-jet branch-local scalar observable is a function of $\eta$ and the longitudinal second derivative
  \[
    Y:=U^\mu U^\nu \nabla_\mu\nabla_\nu\eta.
  \]
\end{enumerate}
If, in addition, the branch is locally $\eta$-complete, then there exists an autonomous reduced equation
\[
  \frac{d^2\eta}{ds^2}=B(\eta)
\]
with
\[
  B(0)=0.
\]
If the branch realizes a closure-restoring return law, then near the exact core one has
\[
  \eta''+\omega^2\eta=O(\eta^2)
\]
with $\omega^2>0$ on the nondegenerate branch. Hence the informal ``Umkehrwalze'' is mathematically realized as a local restoring dynamics for the transported master residual.
\end{theorem}

\begin{proof}
The zeroth-order claim follows from isotropy of the visible defect sector: once $D_{\mu\nu}=\eta g_{\mu\nu}$, all algebraic scalar contractions are polynomials in $\eta$. Branch-local visibility forces the first derivative of any such scalar to be tangent to the same admissible branch direction, so the first jet remains one-dimensional. At second order the only new natural scalar is the longitudinal Hessian component $Y$, and local $\eta$-completeness identifies it as a function of the already visible branch value. The exact core is transport-preserved, so it must be an equilibrium of the reduced equation; closure-restoring return then fixes the linearized sign and yields the stated normal form.
\end{proof}

\begin{definition}[Hybrid reversal probe]
Define the condensed reversal probe by
\[
  \mathfrak S_{\mathrm{rev}}:=(\mathbb P_{\mathrm B},\mathcal L_{\mathrm L}),
\]
where $\mathbb P_{\mathrm B}$ is a hard admissibility/closure projector (the ``Bombe'' side) and $\mathcal L_{\mathrm L}$ is a fine residual functional on the $\eta$-branch (the ``Lorenz'' side).
\end{definition}

\begin{proposition}[Reversal-probe dichotomy]
Let two admissible probe trajectories on the same local residual branch have the same visible master residual value. Then locally only two outcomes remain:
\[
  \Delta\theta\equiv 0 \pmod{2\pi}
  \qquad\text{or}\qquad
  \Delta\theta\equiv \pi \pmod{2\pi}.
\]
Thus the two probes either converge to the same visible branch representative, or they lie on the same transported structure with nontrivial $\mathbb Z_2$ phase class and close only on the spinorial/M\"obius double cover.
\end{proposition}

\begin{proof}
The hard projector removes inadmissible branch continuations, while the fine functional identifies the remaining residual class through the same local variable $\eta$. Equality of the visible residual therefore leaves no generic room for an independent third local phase sector. What remains is either direct coincidence in the visible branch parameterization or the unique nontrivial $\mathbb Z_2$ lift compatible with spinorial sign reversal.
\end{proof}

\subsection*{C.1 First 3.4 extension: closure potential and global branch law}

\begin{definition}[Closure potential and branch energy]
Assume a connected admissible branch interval $I\ni 0$ on which the reduced residual equation is exact, i.e.
\[
  \eta'' = B(\eta) = -\mathcal U'(\eta)
\]
for some $\mathcal U\in C^2(I)$. We call $\mathcal U$ a closure potential for the branch and define the associated branch energy by
\[
  \mathcal E_\eta := \frac12 (\eta')^2 + \mathcal U(\eta).
\]
\end{definition}

\begin{theorem}[Global exact branch law under closure exactness]
\label{thm:global-exact-branch-law}
Assume the hypotheses of Theorem~\ref{thm:eta-completeness-restoring}, and assume in addition branch-local closure exactness in the above sense. Then:
\begin{enumerate}[label=(\roman*)]
  \item $\mathcal E_\eta$ is conserved along every admissible branch trajectory;
  \item the exact compensated core $\eta=0$ is an equilibrium with $\mathcal U'(0)=0$;
  \item if the branch is closure-restoring and nondegenerate, then $\mathcal U''(0)=\omega^2>0$, hence
  \[
    \mathcal U(\eta)=\frac12 \omega^2 \eta^2 + O(\eta^3);
  \]
  \item every admissible turning point of the visible residual branch satisfies
  \[
    \mathcal U(\eta)=\mathcal E_\eta.
  \]
\end{enumerate}
\end{theorem}

\begin{proof}
Differentiating $\mathcal E_\eta$ along a trajectory gives
\[
  \frac{d}{ds}\mathcal E_\eta = \eta'\eta''+\mathcal U'(\eta)\eta' = \eta'(\eta''+\mathcal U'(\eta))=0.
\]
The equilibrium statement is equivalent to $B(0)=0$, already forced by core preservation. The nondegenerate restoring sign from the local normal form gives $B(\eta)=-\omega^2\eta + O(\eta^2)$, hence $\mathcal U'(\eta)=\omega^2\eta + O(\eta^2)$ and therefore the stated Taylor expansion. Turning points are exactly the points where $\eta'=0$, so the conserved-energy identity reduces to $\mathcal U(\eta)=\mathcal E_\eta$.
\end{proof}

\begin{proposition}[Energy-form reversal-probe criterion]
\label{prop:energy-reversal-probe}
Assume the hypotheses of Theorem~\ref{thm:global-exact-branch-law}. Let two admissible probe trajectories be selected by the same hard projector $\mathbb P_{\mathrm B}$ and lie on the same connected regular level set $\mathcal E_\eta = E$. Assume moreover that this energy component carries no monodromy higher than the spinorial $\mathbb Z_2$-lift. Then exactly two global outcomes remain:
\[
\begin{aligned}
&\text{direct convergence on the visible branch}\\[0.2em]
&\qquad \text{or} \qquad
\text{closure only on the spinorial/M\"obius double cover}
\end{aligned}
\]
\end{proposition}

\begin{proof}
The hard projector excludes inadmissible components, while the fine residual functional identifies the same visible residual class through the common potential-energy level. On a connected regular admissible energy component, the remaining ambiguity is topological rather than scalar. By hypothesis no higher monodromy is available, so the only surviving alternatives are trivial monodromy and the unique nontrivial $\mathbb Z_2$ lift.
\end{proof}


\subsection*{C.2 Second 3.4 extension: reversible single-well class}

\begin{theorem}[Parity and single-well restriction of the closure potential]
\label{thm:appendix-single-well-341}
Assume the hypotheses of Theorem~\ref{thm:global-exact-branch-law}. Assume moreover that the admissible residual branch is invariant under the involution $R(\eta)=-\eta$, that the reduced branch law commutes with $R$, and that $\eta=0$ is the unique admissible critical core on the connected interval $I$ with no further admissible critical point of $\mathcal U$ on $I\setminus\{0\}$. Normalize $\mathcal U(0)=0$. Then:
\begin{enumerate}[label=(\roman*)]
  \item $B(-\eta)=-B(\eta)$ and $\mathcal U(-\eta)=\mathcal U(\eta)$ on $I$;
  \item near the exact core,
  \[
    \mathcal U(\eta)=\frac12\omega^2\eta^2+a_4\eta^4+O(\eta^6),
    \qquad
    B(\eta)=-\omega^2\eta-4a_4\eta^3+O(\eta^5);
  \]
  \item $\mathcal U'(\eta)>0$ for every $\eta>0$ in $I$ and $\mathcal U'(\eta)<0$ for every $\eta<0$ in $I$; in particular $\mathcal U(\eta)>0$ for all $\eta\neq0$.
\end{enumerate}
Consequently $\mathcal U$ belongs to the symmetric single-well class on $I$.
\end{theorem}

\begin{proof}
If $\eta(s)$ solves $\eta''=B(\eta)$, branch reversibility and commutation with $R$ imply that $-\eta(s)$ is again an admissible solution of the same reduced law. Hence
\[
  -\eta'' = B(-\eta),
\]
while applying $R$ to the original equation gives $-\eta''=-B(\eta)$, so $B(-\eta)=-B(\eta)$. Since $\mathcal U'(\eta)=-B(\eta)$, the derivative $\mathcal U'$ is odd. With the normalization $\mathcal U(0)=0$, integration yields $\mathcal U(-\eta)=\mathcal U(\eta)$.

Theorem~\ref{thm:global-exact-branch-law} already gives $\mathcal U''(0)=\omega^2>0$. Evenness forces the Taylor series of $\mathcal U$ at the exact core to contain only even powers, so
\[
  \mathcal U(\eta)=\frac12\omega^2\eta^2+a_4\eta^4+O(\eta^6),
\]
and differentiation gives the stated odd expansion of $B$.

Because $\mathcal U''(0)=\omega^2>0$, one has $\mathcal U'(\eta)>0$ for sufficiently small $\eta>0$. If $\mathcal U'$ vanished again at some positive point, continuity would force an additional admissible critical point of $\mathcal U$ on $I\cap(0,\infty)$, excluded by hypothesis. Hence $\mathcal U'(\eta)>0$ on the whole positive half-branch. The negative half-branch follows by oddness of $\mathcal U'$. Therefore $\eta=0$ is the unique well bottom and $\mathcal U(\eta)>0$ for all $\eta\neq0$.
\end{proof}

\begin{proposition}[Symmetric turning-pair criterion]
\label{prop:symmetric-turning-pair-341}
Under the hypotheses of Theorem~\ref{thm:appendix-single-well-341}, let $E>0$ be a regular branch-energy value. Then the turning-point equation
\[
  \mathcal U(\eta)=E
\]
has on $I$ either no solution or exactly the symmetric pair $\eta=\pm\eta_E$ with $\eta_E>0$. In particular every regular admissible energy shell carries at most the scalar ambiguity $\eta\leftrightarrow-\eta$.
\end{proposition}

\begin{proof}
By the preceding theorem, $\mathcal U$ is even and strictly increasing on each half-branch as a function of $|\eta|$. Therefore the equation $\mathcal U(\eta)=E$ can have at most one positive and at most one negative solution, and evenness forces them to occur as a symmetric pair whenever they exist.
\end{proof}

\subsection*{C.3 Third 3.4 extension: canonical quartic family}
\addcontentsline{toc}{subsection}{C.3 Third 3.4 extension: canonical quartic family}

\begin{proposition}[Quartic reduction under first nonlinear visibility]
\label{prop:appendix-342a}
Assume the hypotheses of Theorem~\ref{thm:mainline-342}. Then on a symmetric admissible neighborhood $J$ one has
\[
  \mathcal U(\eta)=\frac12\omega^2\eta^2+\frac{\lambda}{4}\eta^4+R_6(\eta),
  \qquad R_6(\eta)=O(\eta^6),
\]
for some $\lambda\ge 0$.
\end{proposition}

\begin{proof}
By Theorem~\ref{thm:mainline-341}, the potential is even, $C^6$, and has a strict minimum at $0$, so its Taylor expansion on $J$ has the form
\[
  \mathcal U(\eta)=\frac12\omega^2\eta^2+a_4\eta^4+a_6\eta^6+o(\eta^6).
\]
The hypotheses exclude any additional visible scalar or independent local scale at first nonlinear order, so the quartic coefficient is the only admissible datum at that order. Writing $\lambda:=4a_4$ and absorbing the higher terms into $R_6$ gives the claimed form. Since $0$ remains the unique admissible local minimum on $J$, shrinking $J$ if necessary excludes the sign choice that would open a new nearby admissible critical pair; therefore $\lambda\ge 0$ on the retained neighborhood.
\end{proof}

\begin{proposition}[Turning-value parameterization]
\label{prop:appendix-342b}
In the minimal truncation
\[
  \mathcal U_{\min}(\eta)=\frac12\omega^2\eta^2+\frac{\lambda}{4}\eta^4,
  \qquad \lambda\ge 0,
\]
regular energy shells are parameterized by the unique positive turning value $A(E)$ satisfying
\[
  E=\frac12\omega^2A(E)^2+\frac{\lambda}{4}A(E)^4.
\]
If $\lambda>0$, then
\[
  A(E)^2=\frac{-\omega^2+\sqrt{\omega^4+4\lambda E}}{\lambda},
\]
while for $\lambda=0$ one has $A(E)=\sqrt{2E}/\omega$.
\end{proposition}

\begin{proof}
On a regular shell, turning points are exactly the zeros of $\eta'$, so they satisfy $E=\mathcal U_{\min}(\eta)$. By even single-well monotonicity there is exactly one positive solution $A(E)$ and one negative solution $-A(E)$. Solving the quadratic equation in $x=A(E)^2$ gives the explicit formula when $\lambda>0$, and the harmonic formula when $\lambda=0$.
\end{proof}

\begin{proposition}[Appendix C.4: quartic closure coefficient as visible datum]
\label{prop:app-c-34-4}
Under the hypotheses of Theorem~\ref{thm:mainline-343}, the visible quartic branch coefficient $\lambda$ is not an independent nonlinear parameter. It is exactly the transported quartic closure coefficient $\mu c_4$, and every regular probe shell may therefore be parameterized either by its conserved energy $E$ or by its unique positive turning value $A(E)$ without introducing an additional local scalar.
\end{proposition}

\begin{proof}
Theorem~\ref{thm:mainline-343} gives $\lambda=\mu c_4$ with $\mu>0$, so the quartic correction is fixed by the same closure geometry that already determines the quadratic restoring core. Proposition~\ref{prop:appendix-342b} identifies the regular shell by either $E$ or $A(E)$, and Theorem~\ref{thm:mainline-342} excludes any further visible local scalar at quartic order. Hence the quartic visible datum is exactly the transported closure coefficient and does not enlarge the local scalar algebra.
\end{proof}


\begin{proposition}[Appendix C.5: quartic timing drift as probe readout]
\label{prop:app-c-34-5}
Under the hypotheses of Theorem~\ref{thm:mainline-344}, the first nonlinear timing drift of a regular periodic shell is already fixed by the transported quartic closure coefficient. More precisely,
\[
  T(E)=\frac{2\pi}{\omega}\Bigl(1-\frac{3c_4}{4\mu c_2^2}E+O(E^2)\Bigr),
\]
so the branch is locally isochronous at first nonlinear order if and only if $c_4=0$, while $c_4>0$ forces a strictly decreasing period for all sufficiently small regular energies.
\end{proposition}

\begin{proof}
Theorem~\ref{thm:mainline-344} gives the small-energy period expansion together with the closure rewrite $\omega^2=\mu c_2$ and $\lambda=\mu c_4$. Since $\mu>0$ and $c_2>0$, the first nonlinear period correction vanishes exactly when $c_4=0$ and has negative sign whenever $c_4>0$. Thus the timing drift is a direct probe readout of the transported quartic closure coefficient and does not introduce a second nonlinear scalar.
\end{proof}


\begin{proposition}[Appendix C.6: higher even hierarchy and delayed timing defect]
\label{prop:app-c-34-6}
Under the hypotheses of Theorem~\ref{thm:mainline-345}, the first nonzero higher nonlinear timing defect is governed by the first nonvanishing higher even closure coefficient. More precisely, if
\[
  c_4=c_6=\cdots=c_{2m-2}=0,
  \qquad c_{2m}\neq 0,
  \qquad m\ge 3,
\]
then
\[
  T(E)=\frac{2\pi}{\omega}\Bigl(1-\kappa_m(\omega)\,\mu c_{2m}\,E^{m-1}+O(E^m)\Bigr),
  \qquad \kappa_m(\omega)>0.
\]
Hence the branch is isochronous through all lower nonlinear orders, and the first visible higher timing drift appears exactly at order $E^{m-1}$ with sign determined by $c_{2m}$.
\end{proposition}

\begin{proof}
Theorem~\ref{thm:mainline-345} identifies the first visible higher nonlinear coefficient as $\alpha_{2m}=\mu c_{2m}$ and gives the corresponding small-energy period law. Because the intermediate even coefficients vanish, no lower-order contribution to the period expansion is available. Since $\kappa_m(\omega)>0$ and $\mu>0$, the sign of the first nonzero higher timing defect is exactly the sign of $c_{2m}$.
\end{proof}


\begin{proposition}[Appendix C.7: analytic generator form and coefficient transport]
\label{prop:app-c-34-7}
Under the hypotheses of Theorem~\ref{thm:mainline-346}, the local visible branch dynamics is canonically encoded by a single even generator $\Phi$. More precisely,
\[
  \mathcal U(\eta)=\Phi(\eta^2),
  \qquad
  B(\eta)=-2\eta\,\Phi'(\eta^2),
\]
with
\[
  \Phi(s)=\frac12\omega^2 s + \frac{\lambda}{4}s^2 + \sum_{m\ge 3}\frac{\alpha_{2m}}{2m}s^m,
\]
and coefficient transport
\[
  \omega^2=\mu c_2,
  \qquad
  \lambda=\mu c_4,
  \qquad
  \alpha_{2m}=\mu c_{2m}\quad (m\ge 3).
\]
Hence the entire local nonlinear shell hierarchy is fixed by the single generator germ $\Phi$, and every delayed timing defect is a readout of its coefficients rather than a new independent scalar.
\end{proposition}

\begin{proof}
Theorem~\ref{thm:mainline-346} gives the analytic factorization $\mathcal U(\eta)=\Phi(\eta^2)$ together with the coefficient identifications. Expanding $\Phi$ near $s=0$ yields the displayed formula, while differentiating gives the canonical branch law $B(\eta)=-2\eta\Phi'(\eta^2)$. The higher-order timing and turning-value readouts were already tied to the visible coefficients in Theorems~\ref{thm:mainline-344} and \ref{thm:mainline-345}, so the generator form introduces no additional scalar data.
\end{proof}


\begin{proposition}[Appendix C.8: generator-shell classifier]
\label{prop:app-c-34-8}
Under the hypotheses of Theorem~\ref{thm:mainline-347}, every sufficiently small regular admissible shell is completely determined by its generator radius $s(E)=\Phi^{-1}(E)$, and the visible shell ambiguity is exactly the transported sign pair
\[
  \eta=\pm\sqrt{s(E)}.
\]
Hence the hybrid reversal probe has on such a shell only two local transport classes: sign-preserving transport (direct visible convergence) and sign-reversing transport (closure only after the unique spinorial/M\"obius lift).
\end{proposition}

\begin{proof}
Strict monotonicity of $\Phi$ on the small-shell interval gives the inverse shell map $s(E)=\Phi^{-1}(E)$. Since the visible potential is $\mathcal U(\eta)=\Phi(\eta^2)$, the shell equation $\mathcal U(\eta)=E$ is equivalent to $\eta^2=s(E)$, so the shell consists exactly of the two visible representatives $\pm\sqrt{s(E)}$. Any admissible transport restricted to this shell acts on that pair by either preserving sign or reversing it. By the standing no-higher-monodromy hypothesis, these two sign actions are the only local shell transport classes.
\end{proof}


\begin{proposition}[Appendix C.9: shell parity and diagnostic M\"obius-reflected render]
\label{prop:app-c-34-9}
Under the hypotheses of Theorem~\ref{thm:mainline-348}, the local probe outcome on every sufficiently small regular shell is completely determined by the shell parity $\pi(E)\in\{+1,-1\}$. If an admissible shell cycle decomposes into regular segments and transverse reflector crossings, then
\[
  \pi(E)=(-1)^{N_E},
\]
where $N_E$ is the number of transverse reflector crossings in the cycle. Accordingly, even crossing parity yields direct visible convergence, while odd crossing parity yields closure only after the spinorial/M\"obius lift.
\end{proposition}

\begin{proof}
Theorem~\ref{thm:mainline-348} already shows that the shell transport acts on the visible turning pair by a unique sign parity $\pi(E)$. Under the stated decomposition, regular segments preserve sign and every transverse reflector crossing contributes the involution $\eta\mapsto-\eta$. Multiplying these contributions gives $\pi(E)=(-1)^{N_E}$, from which the direct-convergence versus lifted-closure criterion follows immediately.
\end{proof}

\begin{figure}[H]
\centering
\IfFileExists{julia_probe_residual_scales_v3841.pdf}{\includegraphics[width=0.96\textwidth]{julia_probe_residual_scales_v3841.pdf}}{\safegraphic[width=0.96\textwidth]{julia_probe_residual_scales_v3841.png}}
\caption{\textbf{Diagnostic M\"obius-reflected quadratic probe render at the current residual scales.} Comparison of the iterated map $z_{n+1}=1/\overline{(z_n^2+c_\gamma)}$ for the baseline $\gamma=0$ and for the three currently relevant OP~1 residual scales $\Delta_\beta\approx 7.50\times 10^{-5}$, $\Delta_\theta\approx 1.79\times 10^{-4}$, and $\Delta\mu\approx 3.38\times 10^{-4}$. For the present diagnostic render, the parameter family is instantiated by $c_\gamma=-1+i\gamma$, with the unit circle drawn as the M\"obius horizon. The global morphology remains stable across all four renders, while fine boundary filaments and local transition zones respond smoothly under the residual-scale shifts. The figure is included as a diagnostic visualization of structural robustness versus boundary-level response; the earlier marker value $0.0009$ is retained elsewhere only as a historical stress-test scale, not as the current main comparison value.}
\label{fig:app-c-julia-diagnostic}
\end{figure}

Figure~\ref{fig:app-c-julia-diagnostic} is included as a diagnostic probe image rather than as a first-principles derivation. Its role is to visualize a separation between global morphological stability and boundary-level sensitivity: the large-scale visible body remains stable, while the finest branch and transition regions react under the current OP~1 residual scales. The asset is now generated reproducibly by the local Python renderer \texttt{generate\_julia\_probe\_residual\_scales.py}, which reads the displayed values of $\Delta_\beta$, $\Delta_\theta$, and $\Delta\mu$ from the TeX source and produces both a PDF and a PNG version; the document includes the PDF render when available so that the figure can be rebuilt directly from the project material.

\begin{figure}[H]
\centering
\begin{minipage}{0.475\textwidth}
\centering
\safegraphic[width=\linewidth]{gamma0_julia_probe_v3841.png}
\\[-0.5ex]
{\footnotesize $\gamma=0$}
\end{minipage}\hfill
\begin{minipage}{0.47\textwidth}
\centering
\safegraphic[width=0.96\linewidth]{delta_beta_julia_probe_v3841.png}
\\[-0.5ex]
{\footnotesize $\gamma=\Delta_\beta$}
\end{minipage}

\medskip

\begin{minipage}{0.47\textwidth}
\centering
\safegraphic[width=0.95\linewidth]{delta_theta_julia_probe_v3841.png}
\\[-0.5ex]
{\footnotesize $\gamma=\Delta_\theta$}
\end{minipage}\hfill
\begin{minipage}{0.47\textwidth}
\centering
\safegraphic[width=0.95\linewidth]{delta_mu_julia_probe_v3841.png}
\\[-0.5ex]
{\footnotesize $\gamma=\Delta\mu$}
\end{minipage}
\caption{\textbf{Individual diagnostic M\"obius-reflected probe panels.} The four residual-scale renders from Figure~\ref{fig:app-c-julia-diagnostic} shown again as separate panels, so that the boundary-level response at $\gamma=0$, $\Delta_\beta$, $\Delta_\theta$, and $\Delta\mu$ can be inspected without the surrounding comparison frame.}
\label{fig:app-c-julia-panels}
\end{figure}


\subsection*{C.9 Tenth 3.4 extension: tunnel observables and quantitative readout interface}
\addcontentsline{toc}{subsection}{C.9 Tenth 3.4 extension: tunnel observables and quantitative readout interface}

\begin{proposition}[Appendix support for Proposition~\ref{prop:mainline-3410}]
\label{prop:appendix-3410}
Under the hypotheses of Proposition~\ref{prop:mainline-3410}, any admissible quantitative comparison between the diagnostic Julia/probe render and an external readout image may be organized through the tunnel-observable package
\[
  \mathfrak O_{\mathrm{diag}}(\gamma)
  =\bigl(w_{\mathrm{tun}}(\gamma;\tau),\ I_{\mathrm{edge}}(\gamma),\ \Theta_{\mathrm{open}}(\gamma),\ P_{\mathbb Z_2}(E)\bigr),
\]
with no additional local scalar required beyond the generator-shell radius and parity sector. In particular, if two rendering protocols produce the same $s(E)$ and the same parity law on the same thresholded shell family, then any discrepancy among the first three observables is attributable to rendering normalization rather than to a new local structural channel.
\end{proposition}

\begin{proof}
This is a direct restatement of Proposition~\ref{prop:mainline-3410} in appendix language. The tunnel, edge, and opening observables are post-processing functionals of the same thresholded shell geometry, and the parity response is already intrinsic to the shell classifier. Hence matching generator/parity data leaves only normalization choices, not an extra structural scalar.
\end{proof}


\subsection*{C.10 Eleventh 3.4 extension: inverse readout and shell-consistency defect}
\addcontentsline{toc}{subsection}{C.10 Eleventh 3.4 extension: inverse readout and shell-consistency defect}

\begin{theorem}[Appendix support for Theorem~\ref{thm:mainline-3411}]
\label{thm:appendix-3411}
Assume the calibrated tunnel-width law
\[
  w_{\mathrm{tun}}(\gamma;\tau)=W_\tau(s(E))
\]
with strictly monotone $W_\tau$ on a small-shell interval. Then the diagnostic shell data is locally equivalent to the pair
\[
  \bigl(w_{\mathrm{tun}}(\gamma;\tau),P_{\mathbb Z_2}(E)\bigr),
\]
with inverse reconstruction
\[
  s(E)=W_\tau^{-1}(w_{\mathrm{tun}}(\gamma;\tau)),
  \qquad
  \pi(E)=P_{\mathbb Z_2}(E).
\]
Any further shell metric extracted from the same thresholded render therefore factors through these two calibrated quantities.
\end{theorem}

\begin{proof}
Monotonicity of $W_\tau$ gives local invertibility. The shell-parity response already records the transported $\mathbb Z_2$ class. Since Proposition~\ref{prop:mainline-3410} shows that all remaining tunnel/edge/opening metrics are readout-functionals of the same shell geometry, they factor through the recovered pair.
\end{proof}


\subsection*{C.11 Twelfth 3.4 extension: inverse identifiability and stability support}
\addcontentsline{toc}{subsection}{C.11 Twelfth 3.4 extension: inverse identifiability and stability support}

\begin{theorem}[Appendix support for Theorem~\ref{thm:mainline-3412}]
\label{thm:appendix-3412}
Assume the calibrated inverse-readout hypotheses of Theorem~\ref{thm:mainline-3411} and fix a shell point $s_0$ with $W_\tau'(s_0)\neq 0$. Then on a sufficiently small shell neighborhood the calibration law is bi-Lipschitz,
\[
  m_0|s_1-s_2|
  \le |W_\tau(s_1)-W_\tau(s_2)|
  \le M_0|s_1-s_2|,
\]
for suitable constants $m_0,M_0>0$. Hence the calibrated tunnel width determines the shell radius uniquely on each fixed parity branch and does so with inverse condition number bounded by $1/m_0$.
\end{theorem}

\begin{proof}
Since $W_\tau'(s_0)\neq 0$, continuity of the derivative yields a neighborhood on which $W_\tau'$ stays bounded away from zero and from infinity. The mean-value theorem gives the displayed estimate, and the inverse condition-number statement is just the corresponding inverse Lipschitz bound.
\end{proof}


\subsection*{C.13 Fourteenth 3.4 extension: obstruction-shell transition support}
\addcontentsline{toc}{subsection}{C.13 Fourteenth 3.4 extension: obstruction-shell transition support}

\begin{theorem}[Appendix support for Theorem~\ref{thm:mainline-3414}]
\label{thm:appendix-3414}
Let $s_\ast$ be an isolated obstruction shell. If $s_\ast$ is a nondegenerate calibration critical shell, then the calibrated tunnel law has a local fold normal form with a two-sheet inverse exchange. If $s_\ast$ is instead a pure parity shell, then the shell radius continues uniquely and only the transported $\mathbb Z_2$ class jumps. In the mixed case the local transition is the direct product of these two involutions.
\end{theorem}

\begin{proof}
A nondegenerate critical shell of a one-dimensional smooth map is locally a fold, so the inverse consists of exactly two sheets exchanged by the local deck involution. A pure parity shell leaves the scalar calibration regular and therefore does not affect the shell coordinate, but it does flip the transported sign class by definition. Since one operation acts on the shell coordinate and the other on parity, the mixed case is their commuting composition.
\end{proof}

\paragraph{3.4.14 working interpretation.}
The obstruction shells are therefore organized rather than chaotic. Each isolated obstruction contributes only one of two primitive involutions---fold exchange or parity flip---and every local continuation across the inverse-readout atlas is built from these same $\mathbb Z_2$ pieces. This makes the spinorial/M\"obius branch no longer just a visual metaphor at the image layer, but the explicit parity part of the local obstruction monodromy.

\paragraph{3.4.12 working interpretation.}
The inverse-readout block now has a local stability test. Once the calibration slope on a shell interval is known to stay away from zero, tunnel width ceases to be a merely qualitative marker and becomes a controlled coordinate for the shell radius. The remaining local obstruction is discrete---the parity branch---together with the obvious continuous failure mode at critical points of the calibration law.

\subsection*{C.12 Thirteenth 3.4 extension: shell atlas and obstruction support}
\addcontentsline{toc}{subsection}{C.12 Thirteenth 3.4 extension: shell atlas and obstruction support}

\begin{theorem}[Appendix support for Theorem~\ref{thm:mainline-3413}]
\label{thm:appendix-3413}
Let $I$ be a shell interval for the calibrated tunnel-width law $W_\tau$, and let
\[
  \Sigma_{\mathrm{crit}}=\{W_\tau'=0\},
  \qquad
  \Sigma_{\mathrm{par}}=\{\text{parity switch shells}\}.
\]
Then every connected component of $I\setminus(\Sigma_{\mathrm{crit}}\cup\Sigma_{\mathrm{par}})$ is a parity-preserving inverse chart on which $W_\tau$ is one-to-one. On compact subintervals of such a chart, the inverse is bi-Lipschitz and therefore the shell-consistency defect is a well-defined internal diagnostic of model quality rather than of inverse failure.
\end{theorem}

\begin{proof}
Outside $\Sigma_{\mathrm{par}}$ the parity class is fixed. Outside $\Sigma_{\mathrm{crit}}$, the derivative of $W_\tau$ never vanishes and therefore keeps a constant sign on each connected component, so the calibration map is strictly monotone there. Compactness then upgrades the local inverse estimate to a uniform chartwise estimate. The shell-consistency defect is therefore an internal observable on such a chart, and only the two obstruction sets can destroy the chart itself.
\end{proof}

\paragraph{3.4.13 working interpretation.}
The inverse detector-style layer is now stratified. Stable readout happens chart by chart, with sharp obstruction shells where either the calibration slope collapses or the transported parity changes. This gives a clean separation between chart-internal mismatch and genuine breakdown of inverse control.

\paragraph{3.4.11 working interpretation.}
The diagnostic image can therefore be treated as a small inverse instrument. After threshold calibration, tunnel width and shell parity recover the local generator shell and predict the remaining edge/opening readouts. The natural residual mismatch is then measured by the shell-consistency defect $\Delta_{\mathrm{diag}}$, which quantifies how well the image closes back to the same shell/parity model.

\paragraph{3.4.10 working interpretation.}
The image layer is therefore now coupled back into the mathematics in a controlled way: the diagnostic render may be discussed through tunnel width, edge intensification, and opening angle without pretending that visual similarity alone proves a detector interpretation. This keeps the observational hypothesis testable while preserving the structure-first discipline.

\paragraph{3.4.8 working interpretation.}
The shell classifier can therefore be read as a parity law. Once the generator radius $s(E)$ is fixed, the only remaining local ambiguity is whether the admissible shell cycle returns with the same sign or with reversed sign. Under a shell decomposition into regular segments and transverse reflector crossings, this becomes an explicit even/odd criterion. The accompanying M\"obius-reflected diagnostic render is not used as a proof of the residual scalar; it is included only to make the stable-kernel / branch-frontier split visible at the level of the current probe geometry.


\paragraph{3.4.9 working hypothesis.}
The current diagnostic render also motivates a cautious readout hypothesis. Because the plot was not tuned to detector imagery, its visual similarity to some detector-like photon displays should not be dismissed as a mere decorative analogy, but it should likewise not be elevated to an identification claim. Within the standing phase-shift / frame-shift reading, the safer interpretation is that detector images may display readout-relative shadows of the same branch-sensitive residual geometry. This can be tested at the level of tunnel width, opening angle, boundary intensification, and shell-parity response, without assuming that the diagnostic probe already reproduces a full detector event.

\paragraph{3.4.7 working interpretation.}
The analytic generator now also controls the probe decision itself. On every sufficiently small regular shell, the energy first collapses to a unique generator radius $s(E)$, and the visible branch data then reduces to the transported sign pair $\pm\sqrt{s(E)}$. The direct-convergence versus spinorial/M\"obius-lift dichotomy is therefore the $\mathbb Z_2$ shell classifier of the same generator family rather than an additional independent mechanism.

\paragraph{3.4.6 working interpretation.}
The even hierarchy is therefore no longer just a ladder of separate coefficients. Once the transported closure germ is analytic, the visible branch law is packaged by a single local generator $\Phi(\eta^2)$. The quartic shell, the delayed higher timing drifts, and the turning-value corrections then become coordinated readouts of one canonical germ rather than separate nonlinear sectors.

\paragraph{3.4.5 working interpretation.}
The higher nonlinear hierarchy is therefore already strongly constrained: once the quartic channel disappears, the next visible correction cannot re-enter arbitrarily, but only through the first surviving even closure coefficient. Operationally, the probe remains isochronous through all lower orders and only begins to drift at the exact order at which the closure hierarchy itself reopens.

\paragraph{3.4.4 working interpretation.}
The quartic correction is therefore not only structurally present; it is already operationally visible. The first failure of exact isochrony on a regular periodic shell is controlled by the same closure coefficient $c_4$ that generates the quartic branch correction. The probe may therefore read the same admissible shell through energy, through its unique positive turning value, or through its first nonlinear period drift, without opening a new scalar channel.


\paragraph{3.4.3 working interpretation.}
The nonlinear visible coefficient is therefore no longer free: once the exact branch potential is the transported readout of the same closure geometry, the quartic branch coefficient $\lambda$ is simply the positive transport of the closure coefficient $c_4$. The regular shell can then be indexed either by its energy $E$ or by its unique positive turning value $A(E)$, while the nonlinear strength itself is already fixed by the closure side.

\paragraph{3.4.2 working interpretation.}
The first concrete nonlinear model is therefore already visible: once the admissible branch closes to a one-field exact potential and no additional scalar opens at first nonlinear order, the restoring branch law falls into the quartic reversible family. The global probe can then read the same admissible shell either as an energy level or as a unique positive turning value, while the residual ambiguity remains exhausted by the transported $\mathbb Z_2$ pair.

\paragraph{3.4 working interpretation.}
The first move toward 3.4 is therefore not a proliferation of new local lemmas, but a globalization principle: once the reduced branch equation is exact, the restoring branch law becomes a one-dimensional potential dynamics, and the reversal probe becomes an energy-component test for direct closure versus spinorial lift.

\paragraph{Condensed 3.3 status.}
The local 3.3 line is therefore fully captured by four linked claims: a single transported residual scalar $\eta$, a Julia-boundary readout of that same scalar, a restoring branch dynamics for the local reversal profile, and a hybrid reversal probe with the two-way outcome ``direct convergence or spinorial/M\"obius lift''. Relative to the longer 3.3.0 working appendix, this is a structure-preserving compression rather than a deletion: the intermediate defect, jet, and probe lemmas are absorbed into the four main statements above.

